\section{Conclusiones}

En conclusión, la posibilidad de desarrollar dispositivos de medición del porcentaje de \textit{SpO\textsubscript{2}} de manera no invasiva representa una amplia ventaja en el ámbito
clínico ya que permite registrar un signo vital de forma rápida, eficaz y simple. Hay varias posibilidades en el desarrollo de los dispositivos de medición de dicha magnitud, trabajando con amplificadores operacionales,
de transimpedancia, integrados como el MAX305, entre otros. 
La posibilidad de conectar los dispositivos a medidores multiparamétricos hace que sean elementos de uso en una guardia y de conexión rápida al utilizar el conector BC9. Sin embargo, las desventajas radican en la necesidad de conocer qué conexión hay entre 
cada pin, si tienen polaridad o no, y la posibilidad de quemar elementos u obtener mediciones incorrectas si no se coloca correctamente. 

Por otro lado, la realización de un circuito expermiental permitió explorar una implementación distinta de un medidor de saturación de oxígeno, permitiéndonos pensar sobre psibles mejoras y resoución de un problema clínico real. No obstante, la imposibilidad de haber 
detectado la señal pulsátil en el osciloscopio no nos permitió obtener conclusiones específicas sobre la implementación del mismo, y la falta de tiempo dificultó variar magnitudes de los componentes del circuito para ver si se podía solucionar el problema.
Finalmente, la puesta en común con el trabajo de otros equipos logró que se comprendiera lo esperado, que se observaran las curvas pulsátiles y se comprendiera el modo correcto de conexión.